\documentclass[11pt,a4paper]{article}
\usepackage{fontspec}
\setmainfont{DejaVu Sans}
\setsansfont{DejaVu Sans}
\usepackage{polyglossia}
\setdefaultlanguage{russian}
\setotherlanguage{english}
\usepackage{amsmath,amssymb,mathtools}
\usepackage{geometry}
\usepackage{xcolor}
\usepackage{enumitem}
\usepackage{microtype}
\usepackage{fancyhdr}
\usepackage{tcolorbox}
\usepackage{hyperref}
\geometry{left=18mm,right=18mm,top=18mm,bottom=20mm}
\setlength{\parindent}{0pt}
\setlength{\parskip}{4pt}
\definecolor{hseblue}{HTML}{003D7C}
\definecolor{softblue}{HTML}{EEF4FA}
\hypersetup{colorlinks=true,linkcolor=hseblue,urlcolor=hseblue}
\pagestyle{fancy}
\fancyhf{}
\lhead{Методы оптимизации в машинном обучении}
\rhead{Семинар 3}
\cfoot{\thepage}
\newcommand{\R}{\mathbb{R}}
\newcommand{\norm}[1]{\left\lVert #1\right\rVert}
\newcommand{\grad}{\nabla}
\newcounter{problem}
\newenvironment{problem}[1][]{%
  \refstepcounter{problem}%
  \begin{tcolorbox}[colback=white,colframe=hseblue!70!black,boxrule=0.7pt,arc=2mm,left=2.5mm,right=2.5mm,top=2mm,bottom=2mm,title={\textbf{Задача \theproblem.} #1},fonttitle=\normalsize]
}{\end{tcolorbox}}
\begin{document}

{\LARGE\bfseries\color{hseblue} Семинар 3. Квадратичная оптимизация, обусловленность и градиентный спуск}

\vspace{0.4em}
Материал: лекция 3 и первая половина лекции 4. Базовое ядро семинара: задачи 1, 3, 5, 8, 9 и 10. Если не сказано обратное, используется евклидова норма $\norm{\cdot}_2$.

\begin{problem}[Квадратичная функция в матричной форме]
Рассмотрим
\[
f(x,y)=2x^2+2xy+3y^2-4x+2y+5.
\]
Положим $z=(x,y)^\top$.
\begin{enumerate}[label=\alph*)]
\item Представьте функцию в виде
\[
f(z)=\frac12 z^\top Az-b^\top z+c.
\]
Найдите $A,b,c$.
\item Найдите $\grad f(z)$ и $\grad^2f(z)$ непосредственно из координатной записи.
\item Проверьте, что полученные выражения совпадают с формулами
\[
\grad f(z)=Az-b,\qquad \grad^2f(z)=A.
\]
\item Что в этой записи отвечает за форму поверхности, а что --- за положение её центра?
\end{enumerate}
\end{problem}

\begin{problem}[Спектр и положительная определённость]
Пусть
\[
A=
\begin{pmatrix}
3&1\\
1&3
\end{pmatrix}.
\]
\begin{enumerate}[label=\alph*)]
\item Найдите собственные значения и собственные векторы $A$.
\item Определите, является ли $A$ положительно определённой.
\item Используя спектральное разложение, объясните, почему $x^\top Ax>0$ для любого $x\neq0$.
\item Найдите минимальное и максимальное значения отношения Рэлея
\[
R(x)=\frac{x^\top Ax}{x^\top x},\qquad x\neq0.
\]
\end{enumerate}
\end{problem}

\begin{problem}[Точный минимум квадратичной функции]
Пусть
\[
A=
\begin{pmatrix}
2&1\\
1&2
\end{pmatrix},
\qquad
b=
\begin{pmatrix}
1\\
-1
\end{pmatrix},
\]
и
\[
f(x)=\frac12x^\top Ax-b^\top x.
\]
\begin{enumerate}[label=\alph*)]
\item Проверьте, что $A\succ0$.
\item Найдите стационарную точку $x^\star$ из уравнения $Ax^\star=b$.
\item Докажите тождество
\[
f(x)-f(x^\star)=\frac12(x-x^\star)^\top A(x-x^\star).
\]
\item Сделайте вывод о глобальности и единственности найденного минимума.
\end{enumerate}
\end{problem}

\begin{problem}[Собственные направления и геометрия линий уровня]
Рассмотрим
\[
A=
\begin{pmatrix}
5&4\\
4&5
\end{pmatrix},
\qquad
f(x)=\frac12x^\top Ax.
\]
\begin{enumerate}[label=\alph*)]
\item Найдите собственные значения и единичные собственные векторы $A$.
\item В каком собственном направлении функция растёт быстрее при одинаковом перемещении от нуля? Почему?
\item Для линии уровня $f(x)=1$ найдите длины полуосей эллипса в собственных координатах.
\item Под каким углом к координатным осям расположены главные оси этого эллипса?
\end{enumerate}
\end{problem}

\begin{problem}[Число обусловленности и вытянутость эллипса]
Пусть
\[
A=
\begin{pmatrix}
1&0\\
0&25
\end{pmatrix},
\qquad
f(x)=\frac12x^\top Ax.
\]
\begin{enumerate}[label=\alph*)]
\item Найдите $\lambda_{\min}$, $\lambda_{\max}$ и $\kappa_2(A)$.
\item Запишите уравнение линии уровня $f(x)=1$ и найдите её полуоси $a_1,a_2$.
\item Проверьте равенство
\[
\frac{a_{\max}}{a_{\min}}=\sqrt{\kappa_2(A)}.
\]
\item Что изменится в геометрии линий уровня, если заменить $A$ на $7A$? Что произойдёт с $\kappa_2(A)$?
\end{enumerate}
\end{problem}

\begin{problem}[Чувствительность решения линейной системы]
Пусть
\[
A=
\begin{pmatrix}
1&0\\
0&0.02
\end{pmatrix},
\qquad
b=
\begin{pmatrix}
1\\
0
\end{pmatrix},
\qquad
\Delta b=
\begin{pmatrix}
0\\
0.01
\end{pmatrix}.
\]
Рассматриваются системы $Ax=b$ и $A(x+\Delta x)=b+\Delta b$.
\begin{enumerate}[label=\alph*)]
\item Найдите $\kappa_2(A)$, $x$ и $\Delta x$.
\item Вычислите относительное возмущение входных данных
\[
\frac{\norm{\Delta b}_2}{\norm{b}_2}
\]
и относительное изменение решения
\[
\frac{\norm{\Delta x}_2}{\norm{x}_2}.
\]
\item Проверьте оценку
\[
\frac{\norm{\Delta x}_2}{\norm{x}_2}
\le
\kappa_2(A)\frac{\norm{\Delta b}_2}{\norm{b}_2}.
\]
\item Почему в этом примере достигается равенство? Укажите, с какими собственными направлениями совпадают $b$ и $\Delta b$.
\end{enumerate}
\end{problem}

\begin{problem}[Почему возникает градиентный шаг]
Пусть в некоторой точке $x$ известен градиент
\[
g=\grad f(x)=
\begin{pmatrix}
3\\
-4
\end{pmatrix}.
\]
Для $\alpha>0$ рассмотрим вспомогательную функцию
\[
m(h)=g^\top h+\frac1{2\alpha}\norm{h}_2^2.
\]
\begin{enumerate}[label=\alph*)]
\item Найдите $\grad_hm(h)$ и $\grad_h^2m(h)$.
\item Объясните, почему условие $\grad_hm(h^\star)=0$ здесь достаточно для единственного глобального минимума.
\item Найдите $h^\star$ и получите формулу $h^\star=-\alpha\grad f(x)$.
\item Почему минимизируется именно модель $m(h)$, а не исходная функция $f(x+h)$?
\end{enumerate}
\end{problem}

\begin{problem}[Градиентный спуск в собственных координатах]
Пусть
\[
f(x)=\frac12x^\top Ax,\qquad
A=
\begin{pmatrix}
1&0\\
0&4
\end{pmatrix},
\qquad
x_0=
\begin{pmatrix}
1\\
1
\end{pmatrix},
\qquad
\alpha=0.2.
\]
\begin{enumerate}[label=\alph*)]
\item Запишите итерацию градиентного спуска $x_{k+1}=x_k-\alpha Ax_k$.
\item Найдите множители
\[
q_1=1-\alpha\lambda_1,\qquad q_2=1-\alpha\lambda_2.
\]
\item Получите явные формулы для координат $x_{k,1}$ и $x_{k,2}$.
\item Вычислите $x_1,x_2,x_3$.
\item Какая компонента ошибки исчезает быстрее и почему?
\end{enumerate}
\end{problem}

\begin{problem}[Когда градиентный спуск сходится?]
Пусть симметричная положительно определённая матрица $A$ имеет собственные значения
\[
\lambda_1=1,\qquad \lambda_2=5.
\]
\begin{enumerate}[label=\alph*)]
\item Из условий $|1-\alpha\lambda_i|<1$ найдите полный интервал постоянных шагов $\alpha$, при которых градиентный спуск сходится из любой начальной точки.
\item Для $\alpha\in\{0.1,0.25,0.4,0.45\}$ найдите $q_1$ и $q_2$.
\item Для каждого шага опишите поведение ошибки: монотонное затухание, затухающие колебания, пограничный режим или расходимость.
\item Объясните, почему условие
\[
0<\alpha<\frac{2}{\lambda_{\max}}
\]
определяется именно наибольшим собственным значением.
\end{enumerate}
\end{problem}

\begin{problem}[Обусловленность и скорость градиентного спуска]
Пусть
\[
\mu=1,\qquad L=10,\qquad \kappa=\frac L\mu=10,
\]
и используется шаг $\alpha=1/L$.
\begin{enumerate}[label=\alph*)]
\item Для собственного направления с собственным значением $\lambda_i$ запишите
\[
q_i=1-\frac{\lambda_i}{L}.
\]
\item Покажите, что наихудший коэффициент сокращения ошибки достигается при $\lambda_i=\mu$ и равен
\[
q_{\max}=1-\frac{\mu}{L}=1-\frac1\kappa.
\]
\item Найдите численное значение $q_{\max}$.
\item Оцените минимальное целое $k$, гарантирующее
\[
\norm{e_k}_2\le0.01\norm{e_0}_2
\]
из оценки
\[
\norm{e_k}_2\le q_{\max}^k\norm{e_0}_2.
\]
\item Как качественно изменится требуемое число итераций, если $\kappa$ увеличить до $100$?
\end{enumerate}
\end{problem}

\end{document}
